% ===========================================================================
\section{Related work and status of the literature check}\label{sec:prior}
% ===========================================================================

\paragraph{The nearest predecessor}
Leung, Schmidt and Zhang \cite[Lemma~3.11]{LeungSchmidtZhang2023} prove:
if $X=\sum_{i=0}^{p-1}x_i\zeta_p^i$ and $Y=\sum_ix_i$ with
$x_i\in\Z[\zeta_m]$, $p\nmid m$, $p\nmid Y$ and $|X|^2=|Y|^2=pw$,
$p\nmid w$, then $p\mid X\conj Y$. In our notation $Y$ is the
augmentation of $X_\chi$ and $X$ another of its character values, so the
lemma is a statement at the level of character values. The proof of
Theorem~\ref{thm:S} needs two further steps: the conclusion is pushed
from the augmentation to all \emph{coefficients} of $X_\chi$
(Lemma~\ref{lem:column}), and Parseval's identity is applied to a
\emph{single} coset of $C_p$. The column step is a local form of Ma's
lemma \cite{Ma1985}, \cite[Lemma~1.5.1]{Pott1995},
\cite[Result~2.4]{DoDuc2019}; there the input is divisibility of the
character values of all of $G$ by a power of $p$, here it is the norm
equation for the partial evaluation in $\cO[C_p]$.

\paragraph{Other work}
Besides the works cited in Section~\ref{sec:cons}: Brock
\cite{Brock1988} constrains general $\BH(n,h)$; for perfect $p$-ary
sequences and RDS see also Feng \cite{Feng2009}, Liu and Feng
\cite{LiuFeng2015}, and Zhang and Ge \cite{ZhangGe2018} (RDS with
$\gcd(m,n)=1$). Leung and Wang \cite{LeungWang2020} and Leung, Li and Mao
\cite{LeungLiMao2023} prove nonexistence of generalized bent functions
$\Z_2^n\to\Z_m$ by vanishing sums; there $\nu_2(|G|)=n$, so
Theorem~\ref{thm:S} applies only for $n=1$. Ong and Schmidt
\cite{OngSchmidt2026} construct $\BH(\F_q\times\F_q,\operatorname{lcm}
(6,d))$ from Gauss sums; the Sylow subgroups are not cyclic there,
consistent with Remark~\ref{rem:S-sharp}.

\paragraph{Leung--Chue--Zhao}
Leung, Chue and Zhao \cite{LeungChueZhao2025} apply vanishing sums to
cyclic Butson matrices. According to its abstract and to
\cite[Sect.~1]{OngSchmidt2026}, they prove that every $\BH(\Z_p,h)$, $p$
prime, is monomially equivalent to a Fourier matrix (this covers the case
$H=1$ of Theorem~\ref{thm:S}, cf.\ Remark~\ref{rem:H1}); their abstract
also states that they classify the cyclic Butson matrices of order $6$
and give a necessary condition for cyclic Butson matrices of order $2p$.
These results may contain the cases $n=6$ and $n=2p$ of
Corollary~\ref{cor:squarefree}, and the $474$ pairs with $n=2p$
($17$ with $n=6$) among the $2063$ pairs of Do Duc's list in
Section~\ref{sec:cons} may be excluded there.\footnote{Pending: we have
not been able to consult the full text of \cite{LeungChueZhao2025}; the
comparison with it is based on its abstract only and will be completed
before publication.}

\paragraph{Status of the literature check}
We have consulted the full texts (in the versions available to us) of
\cite{DoDuc2019,DoDuc2019b,DucSchmidt2019,LeungSchmidtZhang2023,Lv2017,
Feng2009,Yayla2016,WinterhofYaylaZiegler2014,Arasu2011,
HiranandaniSchlenker2016,LeungSchmidt2005,LeungSchmidt2012,
LeungSchmidt2019,Arxiv260724103,OngSchmidt2026,MaSchmidt1995} and of the arXiv version
of \cite{Schmidt2016survey}, and found no statement of
Theorem~\ref{thm:S}, of its cyclic case, of
Corollary~\ref{cor:squarefree}, or of Theorem~\ref{thm:C} beyond the
cases discussed in Section~\ref{ssec:knownC}. We have \emph{not} been able to consult
the full texts of
\cite{ArasuDeLauneyMa2002,Ma1996,MaNg2009,LeungChueZhao2025,Pott1995,
Schmidt2002LNM,DeLauneyFlannery2011,Mow1995} or of
\cite{Pei1993,Ikeda1999,Feng2001,LiuMaFeng2002,FengLiu2003b}; every
statement we make about these works is taken from the source indicated
at the place where it is used, and the comparison with them is
provisional.
